Este capítulo descreve como o contrato fixado no \refCap{4_Especificacao} se converteu em código Chisel. A ordem de exposição não segue a ordem dos arquivos, mas a ordem das decisões: primeiro a organização dos módulos e a cadeia de dependência entre eles, depois a parametrização, em seguida o arranjo de armazenamento, a separação entre os dois papéis da mesma classe, o controlador e, por fim, a integração com o pipeline e o mecanismo de consistência. Essa ordem corresponde ao percurso efetivamente seguido no desenvolvimento, que foi incremental: cada módulo foi descrito e submetido a teste antes que o seguinte fosse iniciado, de modo que a bancada de um módulo era escrita com a sua especificação ainda em curso.

Vale registrar de antemão um princípio que atravessa o capítulo inteiro. Em Chisel, o texto do programa não é a descrição do circuito; a descrição é o resultado da elaboração desse programa. Toda decisão de implementação discutida a seguir deve ser lida nesse registro: quando se afirma que um valor reside em Reg ou em Wire, a afirmação é sobre o hardware gerado, não sobre a sintaxe.

\section{Estrutura do código e dependência entre módulos}\label{sec:5.1}

O subsistema ocupa o pacote riscv.cache, e a cadeia de dependência entre os seus módulos é linear, o que simplifica tanto a elaboração quanto a verificação isolada de cada um.

{\small
\begin{longtable}{@{}>{\raggedright\arraybackslash}p{(\textwidth - 4\tabcolsep) * \real{0.1732}}>{\raggedright\arraybackslash}p{(\textwidth - 4\tabcolsep) * \real{0.6198}}>{\raggedright\arraybackslash}p{(\textwidth - 4\tabcolsep) * \real{0.2070}}@{}}
\caption[Módulos do subsistema de cache e dependências entre eles]{Módulos do subsistema de cache e dependências entre eles (Fonte: autoria própria)}\label{tab:5.1}\\
\hline
\textbf{Módulo} & \textbf{Responsabilidade} & \textbf{Depende de} \\
\hline
\endfirsthead
\hline
\textbf{Módulo} & \textbf{Responsabilidade} & \textbf{Depende de} \\
\hline
\endhead
\hline
\endlastfoot
CacheParams & geometria, decomposição do endereço e enumeração das políticas & não se aplica \\
CacheIO & interfaces com o pipeline e com a memória, e o agregado de contadores & CacheParams \\
CacheArrays & armazenamento de tag, dado, validade e sujeira & CacheParams \\
CacheLru & interface de reposição e as três políticas & CacheParams \\
Cache & controlador, detecção de acerto e falta, extração e mesclagem & CacheArrays, CacheLru \\
L1Cache & instancia as duas caches e coordena o FENCE.I & Cache \\
CacheMainMemory & memória principal simulada, com arbitragem, latência e erro & CacheParams \\
Pipeline3 & instancia a busca de instrução, o L1Cache e a memória & todos os anteriores \\
\end{longtable}
}

A linearidade não é acidental. CacheParams não conhece nenhum outro módulo e é o único ponto em que as larguras de campo são definidas; CacheArrays e CacheLru guardam estado e não conhecem o controlador; o Cache combina os dois e é o único que interpreta estados; o L1Cache é o único que conhece as duas instâncias ao mesmo tempo; e o Pipeline3 é o único que conhece a hierarquia completa. Nenhuma dependência aponta no sentido inverso, o que permite instanciar e testar qualquer módulo sem construir os que estão acima dele.

A árvore de arquivos reflete essa separação, e a divisão dos testes acompanha a dos pacotes de produção: há um diretório de testes unitários que espelha cada pacote, um diretório específico para a cache, com um teste por aspecto, e um diretório de integração, que exercita o pipeline completo. Essa correspondência entre a estrutura do código e a estrutura dos testes não é formalidade: ela é o que torna localizável o efeito de uma alteração, porque um defeito introduzido em CacheArrays tem um único conjunto de testes como candidato.

\section{A geometria como argumento do gerador}\label{sec:5.2}

A seção~\ref{sec:2.2} sustentou que a vantagem do Chisel sobre as linguagens de descrição tradicionais está em tratar características estruturais do circuito como parâmetros do gerador, e não como constantes repetidas no corpo da descrição. A implementação da geometria é o ponto em que essa afirmação se converte em código.


{\footnotesize
\begin{verbatim}
case class CacheConfig(
    capacityBytes: Int = 8 * 1024,   // capacidade de cada cache
    ways: Int = 4,                   // 1 = mapeamento direto
    lineBytes: Int = 16,             // potência de 2, ao menos 8
    policy: ReplacementPolicy = ReplacementPolicy.Lru) {

  val wordsPerLine = lineBytes / 4
  val numSets     = capacityBytes / lineBytes / ways
  val offsetBits  = log2Ceil(lineBytes)
  val indexBits   = log2Ceil(numSets)
  val tagBits     = 32 - indexBits - offsetBits
  val wayBits     = log2Up(ways)    // ao menos 1 bit, mesmo com 1 via

  require(isPow2(ways) && isPow2(numSets) && numSets >= 2)
}
\end{verbatim}
}


Três decisões merecem comentário. A primeira é que a configuração é um case class, e não um objeto de constantes, o que permite instanciar configurações distintas no mesmo programa de elaboração e compará-las sem alterar o código-fonte. A segunda é que os campos derivados são calculados no construtor, e não recalculados em cada módulo, de modo que a decomposição do endereço da seção~\ref{sec:4.2} exista uma única vez. A terceira é o uso de require, que verifica as restrições da seção~\ref{sec:4.2} em tempo de elaboração: uma configuração inválida falha ao gerar o circuito, e não em simulação, o que é a diferença entre um erro detectado e um circuito silenciosamente incorreto.

Todos os módulos da hierarquia recebem a configuração como parâmetro de construção provido de valor padrão, o que faz com que o código anterior à introdução da parametrização continue compilando sem alteração. Todos os módulos da hierarquia recebem a configuração como parâmetro de construção provido de valor padrão, o que faz com que o código anterior à introdução da parametrização continue compilando sem alteração. O valor padrão é a linha de base do projeto, isto é, 8 KiB por cache, quatro vias, linha de 16 bytes e LRU, e é ele que corresponde aos números da seção~\ref{sec:4.2}.

A política de reposição é escolhida em tempo de elaboração por uma fábrica que devolve módulos distintos a partir de uma interface comum:


{\footnotesize
\begin{verbatim}
object CacheReplacement {
  def apply(cfg: CacheConfig): CacheReplacement = cfg.policy match {
    case ReplacementPolicy.Lru    => Module(new CacheLru(cfg))
    case ReplacementPolicy.Fifo   => Module(new CacheFifo(cfg))
    case ReplacementPolicy.Random => Module(new CacheRandom(cfg))
  }
}
\end{verbatim}
}


A existência dessa fábrica é o que materializa a afirmação da seção~\ref{sec:2.5.4} de que a política de reposição é o parâmetro mais adequado à parametrização: as três alternativas são módulos distintos, com estado interno distinto, e o controlador as consome sem saber qual delas recebeu. A consequência prática é que a comparação entre políticas no \refCap{6_Verificacao} não exige que o controlador seja reescrito, nem que existam três versões do Cache.

\section{O arranjo de armazenamento}\label{sec:5.3}

A linha lida de uma via é descrita como um Bundle, e o conjunto inteiro é devolvido como um Vec desses Bundles:


{\footnotesize
\begin{verbatim}
class CacheLineRead(cfg: CacheConfig = CacheConfig.default) extends Bundle {
  val tag   = UInt(cfg.tagBits.W)                 // tag da linha
  val valid = Bool()                              // linha ocupada?
  val dirty = Bool()                              // difere da memória?
  val data  = Vec(cfg.wordsPerLine, UInt(32.W))   // palavras da linha
}
\end{verbatim}
}


A escrita usa um Bundle distinto, porque a porta de escrita tem semântica mais rica do que a de leitura. A máscara por palavra permite que um preenchimento grave as quatro palavras enquanto um acerto de armazenamento grava apenas a palavra alvo, e o sinal de metadados permite atualizar tag, validade e sujeira sem tocar no dado.


{\footnotesize
\begin{verbatim}
class CacheWritePort(cfg: CacheConfig = CacheConfig.default) extends Bundle {
  val en       = Bool()
  val way      = UInt(cfg.wayBits.W)
  val index    = UInt(cfg.indexBits.W)
  val wordMask = Vec(cfg.wordsPerLine, Bool())    // quais palavras gravar
  val data     = Vec(cfg.wordsPerLine, UInt(32.W))
  val metaEn   = Bool()                           // grava tag/valid/dirty
  val tag      = UInt(cfg.tagBits.W)
  val valid    = Bool()
  val dirty    = Bool()
}
\end{verbatim}
}


No arranjo propriamente dito, a escolha do construto de armazenamento não é uniforme, e a razão está na diferença entre o que precisa ser lido de forma combinacional e o que precisa ser apagado em um ciclo.


{\footnotesize
\begin{verbatim}
private val tagMem  = Seq.fill(cfg.ways)(Mem(numSets, UInt(cfg.tagBits.W)))
private val dataMem = Seq.fill(cfg.ways)(Mem(numSets, Vec(cfg.wordsPerLine, UInt(32.W))))

private val validBits = RegInit(VecInit(Seq.fill(cfg.ways)(
  VecInit(Seq.fill(numSets)(false.B)))))
private val dirtyBits = RegInit(VecInit(Seq.fill(cfg.ways)(
  VecInit(Seq.fill(numSets)(false.B)))))
\end{verbatim}
}


Tag e dado residem em Mem, que tem leitura combinacional e escrita síncrona e admite ser inferida como RAM\glsadd{RAM} distribuída em arranjo de portas programáveis. Validade e sujeira residem em registradores, porque a invalidação exigida pelo FENCE.I precisa apagá-las em todas as vias e todos os conjuntos em um único ciclo, o que uma memória com porta de leitura única não permitiria fazer. Essa separação é a tradução direta da assimetria discutida na seção~\ref{sec:2.5.5} entre drenar, que exige transferir conteúdo, e invalidar, que apenas apaga bits de controle.

A decomposição do endereço é também um Bundle, preenchido por uma função que fatia o endereço conforme a geometria:


{\footnotesize
\begin{verbatim}
class CacheAddress(cfg: CacheConfig) extends Bundle {
  val tag    = UInt(cfg.tagBits.W)
  val index  = UInt(cfg.indexBits.W)
  val offset = UInt(cfg.offsetBits.W)
}
\end{verbatim}
}


Com isso, nenhum módulo da hierarquia manipula índices de campo numéricos: as referências são sempre por nome, e uma mudança de geometria altera as larguras sem alterar os usos.

\section{Um módulo, dois papéis}\label{sec:5.4}

Há uma única classe Cache, e o papel é escolhido por um parâmetro de elaboração. O L1Cache instancia duas cópias:


{\footnotesize
\begin{verbatim}
val instructionCache = Module(new Cache(isInstructionCache = true,  cfg = cfg))
val dataCache        = Module(new Cache(isInstructionCache = false, cfg = cfg))
\end{verbatim}
}


As diferenças entre os dois papéis são todas condicionadas a esse parâmetro, e a tabela abaixo as reúne.

{\small
\begin{longtable}{@{}>{\raggedright\arraybackslash}p{(\textwidth - 4\tabcolsep) * \real{0.2183}}>{\raggedright\arraybackslash}p{(\textwidth - 4\tabcolsep) * \real{0.3600}}>{\raggedright\arraybackslash}p{(\textwidth - 4\tabcolsep) * \real{0.4217}}@{}}
\caption[Diferenças entre a cache de dados e a cache de instruções]{Diferenças entre a cache de dados e a cache de instruções (Fonte: autoria própria)}\label{tab:5.2}\\
\hline
\textbf{Aspecto} & \textbf{DCache} & \textbf{ICache} \\
\hline
\endfirsthead
\hline
\textbf{Aspecto} & \textbf{DCache} & \textbf{ICache} \\
\hline
\endhead
\hline
\endlastfoot
Armazenamento & grava a palavra alvo e marca a linha como suja & o caminho de escrita não é gerado; asserção impede o pedido \\
Leitura & extrai byte, meia-palavra ou palavra, com ou sem extensão de sinal & devolve sempre a palavra inteira de 32 bits \\
Despejo & sim, quando a linha está suja & nunca \\
FENCE.I & drena as linhas sujas & invalida todas as linhas \\
\end{longtable}
}

A diferença de leitura tem justificativa na seção~\ref{sec:4.4}: a cache de instruções entrega a palavra inteira porque a extração do que o estágio de busca precisa depende da codificação da instrução, que a cache não interpreta. A diferença de escrita, por sua vez, decorre de uma exigência de correção, e não de economia: um pedido de escrita dirigido à cache de instruções indicaria erro no estágio consumidor, e o código o trata como violação contratual por asserção, em vez de ignorá-lo em silêncio.

Sobre o bit de sujeira, cabe uma observação. Na cache de instruções a linha nunca fica suja, e isso vale funcionalmente, porque o único ponto que assere o bit é o caminho de armazenamento, que não existe nessa instância, e o preenchimento sempre o grava em zero. Em consequência, a condição que exige despejo é sempre falsa e o estado de despejo é inalcançável nessa cópia. A ausência é também estrutural: o arranjo recebe um parâmetro que condiciona a instanciação dos registradores de sujeira, de modo que os bits não existem no circuito da cache de instruções, e não apenas permanecem constantes em zero. Uma suíte de teste dedicada verifica essa ausência.

\section{O controlador em Reg e Wire}\label{sec:5.5}

O controlador tem oito estados, declarados com Enum e guardados em um único registrador. O mapeamento para os construtos da linguagem obedece a três regras, e é a distinção entre elas que determina o comportamento do circuito.


{\footnotesize
\begin{verbatim}
private val states = Enum(8)
val sIdle           = states(0)   // sem pedido em andamento
val sCompareTag     = states(1)   // compara tags e decide acerto ou falta
val sWriteBack      = states(2)   // grava a vítima suja na memória
val sAllocate       = states(3)   // lê a linha nova da memória
val sFlushScan      = states(4)   // FENCE.I: procura linhas sujas
val sFlushWriteBack = states(5)
val sFlushDone      = states(6)
val sFault          = states(7)   // a memória recusou o acesso

val state = RegInit(sIdle)
\end{verbatim}
}


Reg guarda o que precisa sobreviver entre ciclos: o estado, a via escolhida como vítima e a posição corrente da varredura de drenagem. O registro da vítima merece justificativa própria. Como os arranjos são relidos a cada ciclo e a escolha combinacional depende do conteúdo do conjunto, que muda quando uma linha é gravada, a via da vítima poderia mudar entre o ciclo da falta e o ciclo do despejo. Registrá-la no ciclo da falta é o que preserva a identidade da linha a despejar.

Wire carrega valores do ciclo, e a semântica de última atribuição vence permite um padrão que reduz a probabilidade de erro: as saídas recebem um valor padrão antes da máquina de estados, e cada estado sobrescreve apenas o que efetivamente altera.


{\footnotesize
\begin{verbatim}
io.cpu.hit      := false.B
io.cpu.readData := 0.U
io.cpu.stall    := false.B
io.cpu.fault    := state === sFault
io.mem.request  := false.B
io.mem.write    := false.B

val evHit  = WireDefault(false.B)   // eventos de um ciclo, para os contadores
val evMiss = WireDefault(false.B)
\end{verbatim}
}


Esse padrão tem uma consequência que interessa à verificação: como toda saída tem valor definido em todos os caminhos, não há latch inferido, e a ausência de latch é condição para simular no Verilator, que atualiza apenas na borda de subida do relógio.

O núcleo do controlador é o estado de comparação, e nele se vê que acerto e falta são decididos no mesmo ciclo em que o pedido chega, porque a leitura do arranjo é combinacional:


{\footnotesize
\begin{verbatim}
is(sCompareTag) {
  when(tagHit) {                       // acerto: atende no mesmo ciclo
    io.cpu.hit      := true.B
    io.cpu.readData := loadData
    repl.io.touch   := true.B          // avisa a política de substituição
    state := sIdle
  }.otherwise {                        // falta: guarda a vítima e segue
    io.cpu.stall := true.B
    victimReg    := victimWay
    state := Mux(victimNeedsWriteBack, sWriteBack, sAllocate)
  }
}
\end{verbatim}
}


O fluxo de uma falta é, portanto, uma sequência fechada: comparação, eventual despejo, preenchimento e retorno à comparação. O retorno à comparação, e não ao repouso, é a implementação da decisão justificada na seção~\ref{sec:4.6}: como a escrita no arranjo é síncrona, a linha recém-gravada não seria visível no mesmo ciclo, e o pedido precisa ser reavaliado um ciclo depois para ser atendido como acerto.

O estado de erro completa o conjunto e trata o canal de falha da memória, descrito na seção~\ref{sec:4.5}. Nele o sinal de erro é asserido, a suspensão é liberada e os arranjos permanecem intactos, de modo que nenhuma escrita parcial sobrevive a uma recusa.

\section{A escolha da vítima e as três políticas}\label{sec:5.6}

A escolha da vítima combina duas regras em ordem de precedência, e a primeira delas é o que evita que a política seja consultada desnecessariamente.


{\footnotesize
\begin{verbatim}
val invalidWays = VecInit(setLines.map(l => !l.valid))
val hasInvalid  = invalidWays.asUInt.orR
val victimWay   = Mux(hasInvalid, PriorityEncoder(invalidWays.asUInt), repl.io.victimWay)
\end{verbatim}
}


Se houver via inválida no conjunto, ela é escolhida, e o codificador de prioridade resolve o caso de haver mais de uma. Só quando todas as vias são válidas a política é consultada. A razão é que uma via inválida não contém informação a preservar, de modo que acionar a política nesse caso produziria uma decisão arbitrária sobre um conjunto que ainda não está cheio. O mesmo argumento explica por que a condição que exige despejo é a conjunção entre validade e sujeira: uma vítima inválida não precisa ser gravada na memória, porque não há conteúdo.

As três políticas têm implementações distintas, e o custo de cada uma decorre do que precisa ser mantido por conjunto.

{\small
\begin{longtable}{@{}>{\raggedright\arraybackslash}p{(\textwidth - 6\tabcolsep) * \real{0.1258}}>{\raggedright\arraybackslash}p{(\textwidth - 6\tabcolsep) * \real{0.5215}}>{\raggedright\arraybackslash}p{(\textwidth - 6\tabcolsep) * \real{0.1702}}>{\raggedright\arraybackslash}p{(\textwidth - 6\tabcolsep) * \real{0.1825}}@{}}
\caption[Políticas de reposição e respectivo custo]{Políticas de reposição e respectivo custo (Fonte: autoria própria)}\label{tab:5.3}\\
\hline
\textbf{Política} & \textbf{Estado mantido} & \textbf{Atualização} & \textbf{Custo relativo} \\
\hline
\endfirsthead
\hline
\textbf{Política} & \textbf{Estado mantido} & \textbf{Atualização} & \textbf{Custo relativo} \\
\hline
\endhead
\hline
\endlastfoot
LRU & permutação de posições relativas por conjunto & em todo acerto & o maior \\
FIFO & ponteiro por conjunto & em toda alocação & intermediário \\
Aleatória & registrador de deslocamento em realimentação de 16 bits & em todo ciclo & o menor \\
\end{longtable}
}

A LRU mantém, por conjunto, uma ordenação das vias por recência e a atualiza a cada acerto; é a que exige mais estado, e a que a seção~\ref{sec:2.5.4} escolheu por previsibilidade. A FIFO avança um ponteiro a cada linha alocada, sem considerar uso recente, e o ponteiro é zerado quando a cache de instruções é invalidada, o que é necessário para que o estado da política e o estado do arranjo não se desalinhem. A aleatória usa um registrador de deslocamento sempre em marcha, e a sua vantagem é ser a única das três cujos bits não precisam ser reiniciados na invalidação.

Um detalhe de implementação merece registro, porque conecta esta seção à anterior. A política é avisada do acerto pelo sinal de atualização, e não por um sinal dedicado de alocação. Como o preenchimento termina na comparação, a linha recém-alocada se torna a mais recentemente usada pelo acerto que atende o pedido, e não por uma notificação explícita emitida no preenchimento. O efeito é o descrito na seção~\ref{sec:2.5.4}, e a implementação o obtém sem lógica adicional.

\section{Integração com o pipeline}\label{sec:5.7}

O pipeline tem três estágios, e a cache de dados é acessada no estágio de execução. O ponto que define toda a integração é que o pipeline não lê o sinal de suspensão da cache: ele deriva a espera a partir do sinal de acerto, que é o handshake efetivo. O sinal de suspensão é repassado pelo L1Cache e serve ao contrato de endereço estável e aos testes, mas a decisão de congelar parte de outra premissa.

A escolha de derivar a espera do acerto, e não da suspensão, decorre de uma diferença de semântica entre os dois sinais. A suspensão indica que a cache ainda resolve o pedido, e inclui o ciclo de aceitação; o acerto indica que o dado está disponível naquele ciclo. Como o pipeline precisa congelar até o acerto e liberar no ciclo em que ele ocorre, é o segundo que serve, e o primeiro seria tarde por um ciclo.


{\footnotesize
\begin{verbatim}
// IF: houve pedido à ICache e ainda não houve acerto.
val fetchWait = fetch.io.memoryReadEnable && !l1Cache.io.iCache.hit
// EX: um load ou store pediu à DCache e ainda não houve acerto.
val dataWait  = dataCacheRequest && !l1Cache.io.dCache.hit
// FENCE.I em andamento.
val fenceIWait = fenceIRequest && !l1Cache.io.fenceIDone

cacheWait := fetchWait || dataWait || fenceIWait
stallPipeline := cacheWait || hazardStall
exCanCommit   := !cacheWait && !misalignTrap
\end{verbatim}
}


O efeito do congelamento é total, e não parcial: o contador de programa não avança e os registradores de fronteira entre os estágios mantêm o valor, o que se obtém pela ausência de atribuição sob a condição de espera.


{\footnotesize
\begin{verbatim}
when(cacheWait || hazardStall) {
  pcNext := pcReg                    // PC não avança
}
// IF/ID e ID/EX: sem atribuição, portanto mantêm o conteúdo.
\end{verbatim}
}


O estágio consumidor é a execução, e a condição de commit é o que impede efeitos colaterais durante a espera: com ela falsa, a escrita no banco de registradores, as escritas em registradores de controle e status e a confirmação das instruções atômicas ficam suprimidas até a cache responder.

Duas situações exigem cuidado adicional. A primeira é a espera assimétrica entre as duas caches: se a cache de dados atendeu enquanto a de instruções ainda espera, o resultado precisa ser guardado, sob pena de o acesso ser repetido e, no caso de um armazenamento, gravado duas vezes.


{\footnotesize
\begin{verbatim}
when(cacheWait) {
  when(dataCacheRequest && l1Cache.io.dCache.hit) {
    dataAccessDone := true.B
    dataAccessData := l1Cache.io.dCache.readData
  }
}
\end{verbatim}
}


A segunda é a estabilidade do pedido. Como a seção~\ref{sec:4.4} estabeleceu que o endereço não pode mudar durante a falta, desvios, exceções e a conclusão do FENCE.I só são tomados quando a espera termina:


{\footnotesize
\begin{verbatim}
val takeRedirect = controlRedirect && !cacheWait && !hazardStall
\end{verbatim}
}


A asserção que a cache mantém sobre a estabilidade do endereço funciona, nesse arranjo, como rede de segurança: ela não deve disparar, e o fato de não disparar é o que confirma que a condição de tomada de desvio está correta.

Um último elemento completa a integração. O estágio de busca mantém um buffer com a última palavra recebida, de modo que, ao término de uma falta da cache de instruções, a palavra já esteja disponível e não precise ser pedida novamente. Sem esse buffer, a primeira instrução executada após cada falta custaria um ciclo adicional de busca.

\section{O FENCE.I em duas máquinas}\label{sec:5.8}

A implementação do mecanismo de consistência confirma a decisão de projeto da seção~\ref{sec:4.8}: são duas máquinas distintas, em níveis distintos de hierarquia, e não uma máquina única.

No pipeline, a instrução de sincronização detectada na execução se converte em um pedido ao L1Cache, e durante a espera a busca de instruções é desabilitada:


{\footnotesize
\begin{verbatim}
val fenceIRequest = idEx.valid && !idEx.signals.illegal && idEx.signals.isFenceI
l1Cache.io.fenceI := fenceIRequest
fetch.io.enable   := !fenceIRequest
val fenceIWait     = fenceIRequest && !l1Cache.io.fenceIDone
val fenceIComplete = fenceIRequest && l1Cache.io.fenceIDone
\end{verbatim}
}


A conclusão exige um cuidado que não é evidente. O registrador de fronteira entre busca e decodificação pode conter uma instrução buscada antes da sincronização, e essa instrução pode ser justamente a que foi reescrita. Por isso a conclusão não apenas esvazia as fronteiras: ela refaz a busca a partir do endereço dessa instrução, agora com a memória atualizada e a cache de instruções vazia.


{\footnotesize
\begin{verbatim}
}.elsewhen(fenceIComplete) {
  pcNext := Mux(ifIdValid, ifIdPc, pcReg)
}
\end{verbatim}
}


O buffer do estágio de busca é invalidado no mesmo evento, para que nenhuma palavra antiga seja reaproveitada.

No L1Cache, a máquina tem quatro estados e atua sobre as duas caches:


{\footnotesize
\begin{verbatim}
val fenceIdle :: fenceFlushData :: fenceInvalidateInstruction :: fenceComplete :: Nil = Enum(4)
val fenceState  = RegInit(fenceIdle)
val fenceActive = fenceState =/= fenceIdle

dataCache.io.flushDirty           := fenceState === fenceFlushData
instructionCache.io.invalidateAll := fenceState === fenceInvalidateInstruction
io.fenceIDone                     := fenceState === fenceComplete
\end{verbatim}
}


E, durante a operação, os pedidos das duas caches são filtrados:


{\footnotesize
\begin{verbatim}
instructionCache.io.cpu.request := io.iCache.request && !fenceActive
dataCache.io.cpu.request        := io.dCache.request && !fenceActive
\end{verbatim}
}


O filtro é a implementação da afirmação da seção~\ref{sec:4.8} de que, enquanto a máquina não retorna ao repouso, nenhuma das caches aceita pedido do processador. Ele é necessário porque a drenagem toma a porta de leitura do arranjo: durante a varredura, o índice apresentado ao arranjo é o do conjunto em varredura, e não o do endereço requisitado, de modo que um pedido concorrente não poderia ser atendido ainda que incidisse sobre linha presente.

A drenagem, na cache de dados, percorre os conjuntos e retorna ao mesmo conjunto após cada despejo:


{\footnotesize
\begin{verbatim}
is(sFlushScan) {
  io.cpu.stall := true.B
  when(flushDirtyWays.asUInt.orR) {          // há via suja neste conjunto?
    flushWay  := flushCandidateWay
    flushTag  := flushCandidate.tag
    flushData := flushCandidate.data
    state := sFlushWriteBack
  }.otherwise {
    when(flushSet === (numSets - 1).U) { state := sFlushDone }
    .otherwise { flushSet := flushSet + 1.U }
  }
}
\end{verbatim}
}


O avanço do índice só ocorre quando a varredura de um conjunto não encontra via suja. Como cada passagem trata uma única via e o conjunto é reavaliado na passagem imediatamente posterior, um conjunto com mais de uma via suja é drenado por completo, o que é a implementação do argumento da seção~\ref{sec:4.6}.

Uma diferença entre esta máquina e a de acesso merece registro. No despejo por falta, a via despejada tem o bit de validade asserido e o de sujeira desasserido, porque a linha permanece válida e continua a ser a vítima. Na drenagem, a via é marcada como limpa e também permanece válida, mas o dado não é alterado. Em ambos os casos o que se preserva é a linha, e o que se apaga é a informação de que ela diverge da memória.

Se a memória recusar um despejo durante a drenagem, a operação é abortada e um sinal de erro do FENCE.I indica que restou dado sujo não gravado. Essa indicação é necessária porque, sem ela, a invalidação subsequente apagaria as linhas sujas sem que o conteúdo tivesse chegado à memória, e a perda seria silenciosa, exatamente a consequência que a ordem entre as duas operações existe para evitar.

\section{Ambiente de construção e verificação}\label{sec:5.9}

O projeto usa a versão 3.6.1 do Chisel com Scala 2.13.14 e ChiselTest 0.6.1, conforme justificado na seção~\ref{sec:2.2}. O arquivo de construção declara o plugin de compilador exigido pela linguagem e configura a execução dos testes em processo separado, o que evita que a máquina virtual seja encerrada por falha de simulação. A integração contínua executa a verificação de formatação e a bateria completa de testes a cada alteração, de modo que o estado da linha principal é sempre um estado verificado.

A organização dos testes segue a divisão descrita na seção~\ref{sec:5.1}. O diretório de testes unitários espelha os pacotes de produção; o diretório específico da cache contém um teste por aspecto, a saber, máquina de estados, acerto, política de reposição, arranjo, memória, latência e pedidos simultâneos, canal de erro, contadores e varredura de geometria; e o diretório de integração exercita o pipeline completo, com atenção particular à memória lenta, ao FENCE.I e às instruções comprimidas que cruzam a fronteira entre linhas.

Dois desses testes merecem comentário, porque sustentam afirmações feitas em capítulos anteriores. O teste de varredura de geometria gera as configurações de uma varredura, um eixo por vez com a linha de base nos demais, submete cada uma ao mesmo tipo de sequência de acessos e confere dados e contadores contra um modelo de referência, gravando a taxa de acerto em um arquivo de resultados. É esse teste que materializa a promessa da seção~\ref{sec:2.6} de que o modelo de referência é uma segunda descrição da especificação, e não uma cópia do código. O teste de integração com memória lenta, por sua vez, é o que verifica a afirmação da seção~\ref{sec:5.8} sobre a retomada da busca a partir da instrução reescrita.

A geração de Verilog segue o fluxo descrito na seção~\ref{sec:2.2}: o programa Scala é elaborado, o resultado passa pelo compilador intermediário e o emissor produz o Verilog destinado à simulação e à síntese. O carregamento do programa na memória simulada se dá por uma primitiva de inicialização de memória a partir de arquivo, que se converte em uma diretiva de leitura no Verilog gerado. O arquivo contém uma palavra de 32 bits por linha, em oito dígitos hexadecimais e sem prefixo; a primeira linha corresponde ao endereço zero, a segunda ao endereço quatro, e assim por diante. Programa e dados compartilham o mesmo espaço de endereçamento, e o caminho do arquivo é relativo à raiz do projeto, o que obriga a executar o simulador a partir dela.
